\section{Introduction}
\label{sec:intro}

Hammers~\cite{BlanchetteKaliszykPaulsonUrban2016} connect proof assistants to
automated theorem provers (ATPs): given a goal, they select a set of premises
from the ambient library, translate goal and premises into the ATP's logic ---
usually untyped first-order logic (FOL) with equality --- run external provers,
and reconstruct a found proof inside the proof assistant.  For proof
assistants based on the Calculus of Inductive Constructions (CIC), such as
Rocq (formerly Coq), the translation step must bridge a large gap: CIC is a
dependently typed higher-order logic, while the target is untyped and
first-order.  The CoqHammer translation~\cite{CzajkaKaliszyk2018,Czajka2018}
bridges it by, among other devices, two systematic sources of formula growth:

\begin{enumerate}[label=(\alph*)]
\item \emph{Type guards.}  A quantification $\forall x{:}\tau.\ \phi$ of the
  source logic is translated to
  $\forall x.\ \T(x, \enc{\tau}) \imp \enc{\phi}$, where $\T$ is a binary
  \emph{typing predicate} of the target logic; dually
  $\exists x{:}\tau.\ \phi$ becomes
  $\exists x.\ \T(x,\enc{\tau}) \wedge \enc{\phi}$.  Guards keep the untyped
  provers from instantiating quantifiers with terms of the wrong type, which
  is in general necessary for soundness, but they lengthen every clause,
  enlarge the search space, and are a well-documented performance
  liability~\cite{MengPaulson2008,BlanchetteBohmePopescuSmallbone2016,
  ClaessenLilliestromSmallbone2011}.
\item \emph{Type arguments.}  A polymorphic constant or constructor is
  translated together with its type instantiations: the list constructor
  becomes a ternary symbol in $\mathit{cons}(\alpha, x, l)$.  Type arguments
  inflate every term occurrence and interact badly with term indexing in
  ATPs~\cite{MengPaulson2008,BlanchetteBohmePopescuSmallbone2016}.
\end{enumerate}

Both kinds of annotation are often \emph{redundant}.  If a premise has the
form $\forall x{:}\mathit{nat}.\ \mathit{Even}\,x \imp \phi$, then the guard
$\T(x,\mathit{nat})$ is superfluous \emph{given} the hypothesis atom
$\mathit{Even}(x)$: in the intended interpretation the atom
$\mathit{Even}(x)$ can only hold when $x$ denotes a natural number, because
$\mathit{Even}$ is a predicate \emph{on} natural numbers --- an ill-typed
application $\mathit{Even}\,t$ is not a proposition of the source logic at
all, hence in particular not a provable one.  Similarly, the type argument of
$\mathit{cons}$ is determined by the type of either of its term arguments,
whereas the type argument of $\mathit{nil}$ is determined by nothing that
survives translation and must be kept.  Proof-assistant practice supports the
intuition --- elaboration infers exactly these annotations --- but turning the
intuition into a criterion that is \emph{provably} sound is delicate: the
literature on type encodings for polymorphic first-order and higher-order
logics contains both a series of soundness proofs for guard-based encodings
\cite{ClaessenLilliestromSmallbone2011,BlanchetteBohmePopescuSmallbone2016,
BobotPaskevich2011,CouchotLescuyer2007} and a matching series of concrete
unsoundness examples for naive omissions~\cite{MengPaulson2008,
BlanchetteBohmePopescuSmallbone2016}.  For CIC the situation is harder still:
types are themselves first-class terms, typing is not decidable by sort
inference, and the only published soundness proof for a hammer translation of
(a PTS approximation of) CIC~\cite{Czajka2018} leaves guard omission as a
conjecture and poses the adaptation of monotonicity-based methods as an open
problem.

\paragraph{Contributions.}
This paper gives precise criteria for both omissions and proves them correct
for a substantial fragment of CIC.  Concretely:

\begin{itemize}
\item We define a source calculus \Frag\ --- a polymorphic first-order
  fragment of CIC with parameterized inductive data types, inductive
  predicates, ML-style prenex type quantification, and equality --- together
  with its \emph{ground companion theory} $\Th$, a many-sorted first-order
  theory that captures exactly the reasoning about \Frag-formulas that CIC
  licenses (constructor injectivity, distinctness, case exhaustiveness,
  predicate introduction and inversion), and we make the connection to CIC
  precise (\cref{sec:fragment}).
\item We define the guarded translation $\enc{\cdot}$ of \Frag\ into untyped
  FOL with equality, including the axiom classes emitted by hammer
  translations (typing axioms, unguarded injectivity and discrimination
  axioms, inversion axioms, translated premises) and a class of
  \emph{typehood-completion axioms} $\Comp$ asserting that provable atoms
  carry well-typed arguments.  From an arbitrary many-sorted model $\NN$ of
  $\Th$ we construct an untyped model $\MN$ whose domain contains, besides
  tagged copies of the carriers of $\NN$, a syntactic universe of \emph{junk}
  elements denoting ill-typed applications, and we prove that $\MN$ validates
  the whole translation, completion axioms and unguarded injectivity included
  (\cref{sec:translation}).  This yields a self-contained, model-theoretic
  soundness theorem for the guarded translation: if the translated problem is
  refutable, the goal holds in every model of $\Th$
  (\cref{thm:base-soundness-end}).
\item We formulate the guard-omission criterion (\cref{sec:guards}).  Its
  ingredients are the \emph{typing consequences} $\tc{A}$ of an atom $A$ ---
  guard atoms forced by $A$'s predicate symbol at each argument position ---
  and a pair of mutually recursive syntactic judgments
  $\jtx{x}{u}$/$\jfx{x}{u}$ (``junk-true''/``junk-false'') identifying
  formulas whose truth value is determined under every valuation that
  falsifies the guard $\T(x,u)$.  The main theorem
  (\cref{thm:guard-equivalence}) states that if $\jtx{x}{u}(\phi)$ then
  \[
    \Comp \models (\forall x.\ \T(x,u) \imp \phi) \biimp (\forall x.\ \phi),
  \]
  a logical equivalence relative to the completion axioms; guard omission is
  therefore sound at \emph{any} polarity, in premises and goal alike, and,
  when the completion axioms are emitted, it is also \emph{conservative}: no
  ATP proof is lost (\cref{cor:no-loss}).  We also analyze the existential
  and polarity-restricted cases, where only an implication holds, and prove a
  simultaneous-omission theorem justifying the fixpoint algorithm used in
  practice (\cref{prop:simultaneous}).
\item We formulate the type-argument-erasure criterion (\cref{sec:typeargs}):
  a type argument position $j$ of a function or constructor symbol is
  erasable iff the corresponding type variable occurs in the declared type of
  a retained term argument.  We prove a uniqueness lemma in the style of
  Blanchette et al.'s ``inferable type argument''
  lemma~\cite{BlanchetteBohmePopescuSmallbone2016} --- in our fragment, by an
  elementary injectivity property of type substitution --- and construct from
  $\NN$ an untyped model $\MNe$ of the \emph{erased} translation, in which
  the omitted arguments are recovered uniquely from the retained ones.  Truth
  of the translated goal in $\MNe$ again coincides with truth in $\NN$
  (\cref{lem:erasure-relativization}), which yields soundness of erasure
  (\cref{thm:erasure-soundness}) and of the combination of both optimizations
  (\cref{thm:composition}); the combination requires that predicate symbols
  are never erased, and we show why.
\item We prove the criteria \emph{tight} (\cref{sec:tightness}): dropping an
  uncovered guard, dropping an existential guard conjunct licensed only by
  junk-truth, erasing a result-only type argument, and erasing a predicate's
  type argument each admit a concrete, consistent environment in which the
  transformed translation refutes a goal that fails in some model of $\Th$.
  These four counterexamples are the formal counterparts of unsoundness
  patterns observed in practice~\cite{MengPaulson2008,
  BlanchetteBohmePopescuSmallbone2016,CzajkaKaliszyk2018}.
\end{itemize}

The criteria are designed to be implementable as a purely syntactic
post-processing pass over the first-order formulas produced by an existing
hammer translation, requiring from the source system only the declared types
of the translated constants.  We discuss the relation to guard-based and
monotonicity-based type encodings, to Mizar's cluster registrations, and to
the proof-theoretic soundness framework of~\cite{Czajka2018} in
\cref{sec:related}, together with the limitations of the fragment
(higher-order features, dependent function types, conversion) and the
corresponding open problems.
